% ============================================================
% SEÇÃO 1 – MODELO DESCRITIVO
% ============================================================

\section{MODELO DESCRITIVO}

\subsection{INTRODUÇÃO}

\subsubsection{Objetivos do projeto}

LifeCity tem como objetivo geral desenvolver um aplicativo mobile que facilite a
comunicação entre moradores e poder público, promovendo maior transparência,
participação social e agilidade na identificação e resolução de problemas urbanos.
A proposta central é criar um canal digital acessível e intuitivo que permita ao
cidadão acompanhar informações relevantes sobre sua região, registrar reclamações,
solicitar melhorias e visualizar o andamento dessas demandas de forma organizada
e transparente para as autoridades responsáveis.

Como objetivos específicos, o projeto busca oferecer aos usuários um ambiente em
que possam publicar relatos sobre problemas urbanos, como buracos nas vias, falta
de iluminação, descarte irregular de lixo e outras situações que afetam diretamente
a qualidade de vida da comunidade. Além disso, pretende-se disponibilizar um feed
interativo e geolocalizado, semelhante ao de redes sociais, que permita visualizar
demandas, atualizações, iniciativas da região e também eventos locais como festas,
atividades esportivas, encontros culturais e grupos organizados, incentivando o
engajamento entre os moradores.

Outro objetivo importante é disponibilizar ao poder público uma ferramenta que
concentre denúncias, solicitações e informações sobre o território, facilitando a
priorização, o planejamento e o acompanhamento das ações. Ao consolidar esses
dados, o LifeCity contribui para tomadas de decisão mais eficientes e baseadas em
evidências. Por fim, o projeto busca fortalecer o senso de comunidade, aproximando
indivíduos que compartilham o mesmo espaço urbano e estimulando a colaboração
em prol de melhorias coletivas e de uma vida urbana mais ativa e conectada.

\subsubsection{Cenário atual do mercado}

Atualmente, o mercado de soluções digitais voltadas à gestão urbana e à
participação cidadã tem crescido significativamente, impulsionado pelo aumento do
uso de smartphones e pela necessidade de modernização dos serviços públicos.
Diversas cidades, especialmente as de médio e grande porte, vivem um processo
de transformação digital, buscando ferramentas que melhorem a comunicação entre
governo e população, facilitem o monitoramento de demandas locais e reforcem a
transparência administrativa. No entanto, muitas dessas soluções ainda são
fragmentadas, pouco acessíveis ou não integram de forma eficiente os diferentes
atores envolvidos no cuidado com a cidade. Esse contexto abre espaço para
plataformas inovadoras que centralizam informações, potencializam a voz dos
moradores e otimizam a atuação do poder público.

\paragraph{Cenário atual do usuário}

O cidadão da região sente dificuldade de acessar informações regionais de maneira
organizada e confiável. Grande parte dos cidadãos depende de canais dispersos
como grupos de redes sociais, atendimento presencial em repartições públicas ou
números de telefone específicos para registrar reclamações ou solicitar melhorias.
Esses meios, além de pouco padronizados, muitas vezes não garantem retorno,
deixando o usuário sem visibilidade sobre o andamento do seu pedido. Além disso,
existe uma sensação crescente de distanciamento entre população e gestão
municipal, o que reforça a necessidade de ferramentas que promovam maior
participação cidadã, estimulando o engajamento comunitário e dando mais
autonomia para que o morador acompanhe a situação de sua região em tempo real.

\paragraph{Cenário atual dos poderes públicos}

Do ponto de vista do poder público, o cenário atual apresenta desafios relacionados
à gestão eficiente das demandas da população e à comunicação direta com os
cidadãos. Os órgãos municipais frequentemente lidam com um alto volume de
reclamações provenientes de diferentes canais, dificultando a organização, a
priorização e o acompanhamento das solicitações. A falta de um sistema
centralizado compromete a transparência, torna o processo mais lento e pode gerar
insatisfação por parte da população. Ao mesmo tempo, administrações municipais
buscam soluções tecnológicas que permitam monitorar problemas urbanos com
maior agilidade, gerar relatórios precisos e tomar decisões baseadas em dados.
Dessa forma, há um espaço claro no mercado para ferramentas que conectem, de
forma estruturada, governo e moradores, promovendo uma gestão mais eficiente e
aproximando o poder público das necessidades reais da comunidade.

\subsection{MOTIVAÇÃO}

A motivação para o desenvolvimento do LifeCity surge da percepção de que, mesmo
com os avanços tecnológicos e a presença constante dos smartphones no cotidiano
das pessoas, a comunicação entre moradores e poder público ainda é marcada por
falhas, lentidão e pouca transparência. Muitos cidadãos não sabem a quem recorrer
quando enfrentam problemas urbanos como iluminação precária, acúmulo de lixo ou
vias danificadas e quando conseguem registrar uma reclamação, muitas vezes não
recebem retorno nem têm meios claros para acompanhar o andamento da solução.
Essa distância entre população e gestão municipal alimenta frustrações, reduz o
senso de pertencimento e enfraquece a participação social.

Além disso, grande parte das discussões sobre a cidade acontece em canais
informais, como grupos de redes sociais, onde relatos de problemas e até
divulgações de eventos se perdem entre mensagens, sem alcançar de forma
organizada os responsáveis ou o público mais amplo. Essa desorganização prejudica
tanto o morador, que não vê resultado nem visibilidade para suas demandas, quanto
o poder público, que deixa de aproveitar um fluxo estruturado de informações para
identificar padrões, priorizar ações e otimizar recursos.

O LifeCity nasce justamente para organizar esse fluxo. Ao oferecer um ambiente
digital em que qualquer cidadão possa registrar reclamações, sugestões e
observações de forma simples, geolocalizada e transparente, o aplicativo estimula o
engajamento cívico e fortalece o sentimento de responsabilidade coletiva.
Paralelamente, ao permitir o cadastro e o acompanhamento de eventos locais como
festas, atividades esportivas, encontros culturais e iniciativas comunitárias, o
LifeCity amplia o olhar para além dos problemas, promovendo também aquilo que
fortalece a convivência e a identidade urbana.

Do ponto de vista do poder público, o projeto busca disponibilizar uma ferramenta
prática de monitoramento urbano, concentrando denúncias, solicitações e
informações sobre o território em um único ambiente. Com isso, decisões podem
ser tomadas com base em dados reais e atualizados, aumentando a eficiência na
resposta às demandas da população.

Por fim, o LifeCity se fundamenta em um compromisso com a inovação social:
utilizar a tecnologia para transformar a relação entre cidadãos e cidade, fortalecendo
a comunicação, ampliando o acesso à informação e contribuindo para um espaço
urbano mais organizado, participativo, conectado e vivo, tanto na resolução de
problemas quanto na promoção de eventos e iniciativas comunitárias.

\subsubsection{Respostas das pesquisas feitas com usuários}

Através do questionário realizado no Google Forms, as informações de 16 respostas
foram retiradas para validação do aplicativo. Nesta pesquisa, 68\% dos usuários
estão entre a faixa etária de 22 a 38 anos.

\begin{figure}[H]
  \centering
  \includegraphics[width=0.7\textwidth]{images/grafico1_faixa_etaria.png}
  \caption{Gráfico referente à faixa etária dos usuários}
\end{figure}

Dentre os respondentes, aproximadamente 50\% residem na cidade de Campinas.

\begin{figure}[H]
  \centering
  \includegraphics[width=0.7\textwidth]{images/grafico2_cidade.png}
  \caption{Gráfico referente à cidade de residência dos respondentes}
\end{figure}

100\% dos respondentes afirmam ter presenciado problemas urbanos em sua região
recentemente.

\begin{figure}[H]
  \centering
  \includegraphics[width=0.6\textwidth]{images/grafico3_problemas.png}
  \caption{Confirmação de problemas presenciados na região}
\end{figure}

Os remetentes puderam relatar os problemas urbanos que presenciaram; dentre
esses, a maioria se trata de má infraestrutura das vias públicas e falta de segurança.

81\% dos remetentes da pesquisa afirmam ver muita utilidade no aplicativo.

\begin{figure}[H]
  \centering
  \includegraphics[width=0.6\textwidth]{images/grafico5_utilidade.png}
  \caption{Confirmação de necessidade do aplicativo}
\end{figure}

As funcionalidades do aplicativo foram vistas com importância de forma diversa.

\begin{figure}[H]
  \centering
  \includegraphics[width=0.8\textwidth]{images/grafico6_funcionalidades.png}
  \caption{Funcionalidades mais importantes para o usuário}
\end{figure}

As funcionalidades relacionadas a eventos e cultura também tiveram uma
porcentagem alta de interesse por parte dos usuários da pesquisa. 62\% dos
respondentes afirmam ver necessidade das funcionalidades relativas a atividades
culturais. Boa parte dos respondentes demonstram insatisfação em relação à forma
como a divulgação de atividades culturais ocorre atualmente.

\subsubsection{Oportunidades Encontradas}

O desenvolvimento do LifeCity apresenta oportunidades significativas tanto do ponto
de vista social quanto tecnológico. Ao criar uma plataforma que centraliza
informações, denúncias e interações entre moradores e gestores públicos, o projeto
abre espaço para o fortalecimento da participação cidadã e para a construção de
uma gestão urbana mais eficiente e transparente. Além disso, o aplicativo
representa uma oportunidade de inovação no mercado de soluções cívicas,
ampliando a digitalização de serviços públicos e oferecendo novas possibilidades de
análise de dados territoriais para tomada de decisões. Do ponto de vista acadêmico
e profissional, o LifeCity também permite explorar tecnologias modernas, como
Flutter e Node.js, consolidando conhecimentos em desenvolvimento mobile,
integração com serviços em nuvem e criação de interfaces centradas na experiência
do usuário. Essas oportunidades reforçam a relevância do projeto e seu potencial de
impacto positivo na sociedade.

\subsection{CÍRCULO DOURADO}

\subsubsection{Por quê?}

Porque acreditamos que todo cidadão merece ser ouvido e ter uma cidade mais
unida, organizada, transparente e conectada com suas reais necessidades.

\subsubsection{Como?}

Fazemos isso criando uma plataforma intuitiva que aproxima moradores e poder
público, centralizando reclamações, informações e demandas de forma clara e
acessível.

\subsubsection{O quê?}

Entregamos um aplicativo mobile que funciona como um feed colaborativo da
cidade, permitindo que usuários publiquem problemas urbanos e acompanhem sua
resolução.

\subsection{REGRAS DE NEGÓCIO}

\begin{enumerate}
  \item \textbf{Informações verdadeiras no cadastro:} O usuário deve fornecer dados
        reais para garantir a autenticidade do perfil e permitir comunicação confiável.
  \item \textbf{Uso adequado da plataforma:} O aplicativo deve ser utilizado apenas
        para relatar problemas urbanos, solicitar melhorias e compartilhar informações
        relevantes sobre a cidade.
  \item \textbf{Postagens responsáveis:} As publicações devem conter informações
        verdadeiras, fotos reais e descrições claras, sem exageros, manipulações ou
        denúncias falsas.
  \item \textbf{Respeito e cordialidade:} É proibido publicar conteúdo ofensivo,
        discriminatório, difamatório, inadequado ou que incentive violência ou preconceito.
  \item \textbf{Proibição de uso indevido:} Não é permitido utilizar o aplicativo para
        autopromoção, campanhas políticas, publicidade não autorizada, spam ou qualquer
        finalidade alheia ao objetivo da plataforma.
  \item \textbf{Proteção à privacidade:} É proibido divulgar dados pessoais de
        terceiros, imagens de indivíduos sem autorização ou qualquer informação que
        viole a privacidade de outras pessoas.
  \item \textbf{Conteúdos relacionados à cidade:} As publicações devem estar
        estritamente relacionadas a problemas, avisos ou situações do espaço urbano e
        da comunidade local.
  \item \textbf{Acompanhamento responsável das demandas:} Os usuários devem
        utilizar os recursos do aplicativo de maneira consciente, acompanhando o
        andamento das solicitações e interagindo de forma construtiva.
  \item \textbf{Consequências para uso indevido:} Reclamações falsas, comportamento
        inadequado ou violação das regras podem resultar em advertência, suspensão ou
        exclusão da conta.
\end{enumerate}

\subsection{DESCRIÇÃO DOS REQUISITOS FUNCIONAIS}

\begin{longtable}{|p{1.5cm}|p{3cm}|p{6cm}|p{3.5cm}|}
  \hline
  \textbf{ID} & \textbf{Nome} & \textbf{Descrição} & \textbf{Informação/Dados} \\
  \hline
  \endfirsthead
  \hline
  \textbf{ID} & \textbf{Nome} & \textbf{Descrição} & \textbf{Informação/Dados} \\
  \hline
  \endhead
  RF001 & Cadastro & Criação de conta com informações do usuário & Nome, e-mail, senha, CPF, telefone \\
  \hline
  RF002 & Login & Acesso ao sistema via autenticação JWT & E-mail, senha \\
  \hline
  RF003 & Registrar reclamação & Registro de problema urbano com texto, categoria e localização & Descrição, categoria, localização, foto (opcional) \\
  \hline
  RF004 & Editar/Excluir reclamação & Permite ajustar e remover suas reclamações & ID da reclamação \\
  \hline
  RF005 & Registrar evento & Registro de evento urbano com texto e localização & Descrição, categoria, localização, data/hora \\
  \hline
  RF006 & Editar/Excluir evento & Permite ajustar e remover seus eventos & ID do evento \\
  \hline
  RF007 & Visualizar mapa & Mapa interativo com pins de eventos e reclamações & Coordenadas, tipo \\
  \hline
  RF008 & Filtrar por categoria & Filtra o mapa por categoria (infraestrutura, segurança, limpeza, trânsito, festas, esportes, educação) & Categoria \\
  \hline
  RF009 & Perfil do usuário & Exibe e edita nome, foto e dados básicos & Nome, foto, telefone, CPF \\
  \hline
  RF010 & Configurações & Permite alternar tema (claro/escuro), alterar senha e alterar telefone & Preferências do usuário \\
  \hline
  RF011 & Notificações & Exibe notificações de curtidas, comentários, pedidos de amizade, conquistas e missões com badge de não lidas & Dados da notificação \\
  \hline
  RF012 & Sistema de XP e Níveis & Calcula XP do usuário por ações cívicas (criar reclamação, receber curtida, comentário) e exibe nível e barra de progresso no perfil & XP acumulado, nível atual \\
  \hline
  RF013 & Conquistas & Desbloqueia conquistas automaticamente ao atingir marcos (ex.: 1\textsuperscript{a} reclamação, 10 curtidas recebidas); usuário pode destacar até 3 no perfil & ID da conquista, tipo de gatilho \\
  \hline
  RF014 & Missões diárias e semanais & Sorteia missões do pool de templates para cada usuário diariamente/semanalmente e exibe progresso em tempo real & Tipo de frequência, meta, progresso atual \\
  \hline
  RF015 & Equipes & Permite criar equipes permanentes (mín. 2 / máx. 7 membros), convidar amigos e receber bônus de XP semanal coletivo & Nome da equipe, membros, XP total \\
  \hline
  RF016 & Social — amizades e curtidas & Permite adicionar amigos, curtir e comentar reclamações e visualizar o perfil público de outros usuários & ID do usuário-alvo \\
  \hline
\end{longtable}

\subsection{DESCRIÇÃO DOS REQUISITOS NÃO-FUNCIONAIS}

\begin{longtable}{|p{4cm}|p{10cm}|}
  \hline
  \textbf{ID} & \textbf{Descrição} \\
  \hline
  \endfirsthead
  \hline
  \textbf{ID} & \textbf{Descrição} \\
  \hline
  \endhead
  RNF001 – Tempo de resposta & O sistema deve exibir o feed e carregar dados em até 2 segundos em condições normais de rede. \\
  \hline
  RNF002 – Estabilidade & O aplicativo deve suportar uso contínuo sem travamentos. \\
  \hline
  RNF003 – Eficiência em Consultas & Operações de leitura no banco de dados devem ser otimizadas para reduzir carga excessiva. \\
  \hline
  RNF004 – Experiência do Usuário (UX) & A interface deve ser intuitiva, consistente e fácil de utilizar, seguindo boas práticas de design. \\
  \hline
  RNF005 – Acessibilidade & O aplicativo deve oferecer fonte legível, contraste adequado e elementos interativos acessíveis. \\
  \hline
  RNF006 – Navegação Simplificada & O usuário deve concluir ações principais em no máximo 3 toques. \\
  \hline
  RNF007 – Suporte Multiplataforma & O aplicativo deve ser compatível com Android e iOS via Flutter. \\
  \hline
  RNF008 – Adaptabilidade & As telas devem se ajustar a diferentes tamanhos de dispositivos móveis. \\
  \hline
  RNF009 – Autenticação Segura & A autenticação deve utilizar tokens JWT de acesso e refresh, com renovação automática e proteção de credenciais. \\
  \hline
  RNF010 – Criptografia de Dados & Dados sensíveis devem ser armazenados de forma segura conforme boas práticas. \\
  \hline
  RNF011 – LGPD & O sistema deve garantir privacidade e tratamento adequado dos dados pessoais, respeitando a Lei Geral de Proteção de Dados. \\
  \hline
  RNF012 – Controle de Acesso & Cada usuário acessa apenas suas informações privadas. \\
  \hline
  RNF013 – Disponibilidade & O serviço backend deve permanecer disponível 24/7, salvo períodos de manutenção. \\
  \hline
  RNF014 – Tolerância a Falhas & O sistema deve tratar erros de conexão exibindo mensagens adequadas ao usuário. \\
  \hline
  RNF015 – Integridade dos Dados & Toda informação armazenada deve manter consistência, evitando duplicidade e perda de dados. \\
  \hline
  RNF016 – Código Modular & O código em Flutter/Node.js deve ser organizado de forma modular, facilitando manutenção e expansão futura. \\
  \hline
  RNF017 – Versionamento & Todas as atualizações devem ser controladas via Git. \\
  \hline
  RNF018 – Suporte a Temas & O aplicativo deve oferecer modo claro e modo escuro, com alternância controlada pelo usuário nas configurações. \\
  \hline
\end{longtable}

\subsection{TECNOLOGIAS E LINGUAGENS DE PROGRAMAÇÃO UTILIZADAS}

O desenvolvimento do LifeCity utilizou um conjunto de tecnologias modernas que
apoiam desde a prototipação inicial até a implementação do aplicativo e a
colaboração em equipe. A seguir são apresentadas as principais ferramentas,
linguagens e plataformas adotadas no projeto.

\subsubsection{Linguagem de Programação e Framework}

\textbf{Flutter / Dart}
\begin{itemize}
  \item Framework utilizado para o desenvolvimento do aplicativo mobile.
  \item Permite criar interfaces nativas para Android e iOS a partir de um único
        código-fonte.
  \item Escolhido pela alta produtividade, ampla comunidade e disponibilidade de
        plugins.
  \item Bibliotecas principais utilizadas: \texttt{flutter\_map} (mapa
        OpenStreetMap), \texttt{dio} (HTTP), \texttt{provider} (gerenciamento de
        estado), \texttt{google\_fonts}.
\end{itemize}

\textbf{Node.js / Express}
\begin{itemize}
  \item Utilizado como backend principal da aplicação, responsável por toda a
        camada de negócio e autenticação.
  \item Implementa API RESTful consumida pelo aplicativo Flutter via HTTP.
  \item Gerencia autenticação com tokens JWT (access token + refresh token),
        incluindo renovação automática de sessão.
  \item Fornece endpoints para usuários, reclamações e eventos.
\end{itemize}

\subsubsection{Banco de Dados e Armazenamento}

\textbf{PostgreSQL}
\begin{itemize}
  \item Banco de dados relacional utilizado para armazenar usuários, reclamações
        e eventos.
  \item Gerenciado via Supabase.
\end{itemize}

\textbf{Supabase}
\begin{itemize}
  \item Plataforma open-source utilizada exclusivamente para:
  \begin{itemize}
    \item Banco de dados PostgreSQL gerenciado;
    \item Armazenamento de arquivos de mídia (fotos de perfil e reclamações).
  \end{itemize}
  \item A autenticação \textbf{não} utiliza o serviço nativo do Supabase; é
        realizada inteiramente pelo backend Node.js com JWT próprio.
\end{itemize}

\subsubsection{Autenticação}

\textbf{JWT (JSON Web Tokens)}
\begin{itemize}
  \item Implementação própria no backend Node.js.
  \item O fluxo utiliza \textit{access token} (curta duração) e
        \textit{refresh token} (longa duração).
  \item O aplicativo Flutter renova o access token automaticamente ao receber
        resposta 403, sem necessidade de novo login pelo usuário.
\end{itemize}

\subsubsection{Prototipação e Design}

\textbf{Figma}
\begin{itemize}
  \item Ferramenta utilizada para wireframes, protótipos de baixa e alta
        fidelidade e documentação visual.
  \item Facilitou a construção do fluxo de telas e o alinhamento visual entre
        os membros da equipe.
\end{itemize}

\subsubsection{Comunicação e Organização}

\textbf{Discord}
\begin{itemize}
  \item Utilizado como meio principal de comunicação, reuniões e tomada de
        decisões durante o desenvolvimento.
\end{itemize}

\textbf{Git e GitHub}
\begin{itemize}
  \item Utilizados para versionamento do código-fonte.
  \item Permitem controle de versões, trabalho colaborativo e rastreamento de
        histórico de desenvolvimento.
\end{itemize}
